\chapter{Simulation Stack Development}
\label{ch:sim-stack}

% Guideline recommends 10-20 pages per methodology chapter.
% TODO: add figures - a URDF joint-tree diagram, a screenshot of the
% Gazebo model, and an odometry-drift plot from Sec. 3.6 would all
% strengthen this chapter with minimal extra writing.

This chapter documents the development of the simulation-side
navigation stack that Chapters~\ref{ch:hardware-bridge}
through~\ref{ch:nav-integration} later migrate onto physical
hardware. The stack went through two full drive-architecture
migrations before reaching the form used for hardware bridging
(Sections~\ref{sec:baseline-manual}--\ref{sec:second-migration}), and
along the way surfaced a middleware-level parser defect
(Section~\ref{sec:colon-bug}) and a four-part vehicle-dynamics bug
pattern that recurred, largely unchanged, across two sibling packages
(Section~\ref{sec:not-moving-bug}). The chapter closes with a
behaviour-tree fix specific to Ackermann-steered vehicles
(Section~\ref{sec:bt-fix}).

\section{Baseline: Manual Bicycle-Model Inversion}
\label{sec:baseline-manual}

The navigation stack's earliest drive architecture used two
independent \texttt{ros2\_control} \texttt{JointGroup} controllers
(\texttt{drive\_controller} for wheel velocity,
\texttt{steering\_controller} for hub position) driven by a
hand-written node, \texttt{cmd\_vel\_to\_drive.py}, that manually
inverted the bicycle-model relationship of
Equations~\ref{eq:bicycle-x}--\ref{eq:bicycle-theta} to compute a
per-wheel velocity and steering-hub angle from an incoming
\texttt{geometry\_msgs/Twist} on \texttt{/cmd\_vel}. This
architecture worked, but duplicated the Ackermann inverse-kinematics
logic that a purpose-built controller would otherwise provide, and --
more importantly for a navigation stack meant to reach physical
hardware -- could not support Nav2's recovery behaviours that require
reversing, since the manual node had no reverse-driving logic of its
own.

\section{First Migration: \texttt{ros2\_control}'s \texttt{AckermannSteeringController}}
\label{sec:first-migration}

The first architectural migration replaced both custom
\texttt{JointGroup} controllers and \texttt{cmd\_vel\_to\_drive.py}
with \texttt{ros2\_control}'s own
\texttt{ackermann\_steering\_controller} / \texttt{AckermannSteeringController}
plugin, configured via \texttt{config/qcar\_controllers.yaml}. This
controller performs the same Ackermann inverse kinematics internally,
and additionally publishes its own wheel-based odometry -- which
this project deliberately adopted as the vehicle's \texttt{/odom}
source in place of Gazebo's \texttt{gazebo\_ros\_p3d} ground-truth
plugin (the latter was kept running, but remapped to
\texttt{/ground\_truth/odom} for debugging comparison only). This was
a considered choice, not a default: ground-truth odometry is
trivially available in simulation but has no equivalent on the
physical robot, whose odometry can only ever be a real,
drift-prone dead-reckoning estimate (Section~\ref{sec:onboard-bridge-odom}).
Developing and tuning the navigation stack against the same
\emph{kind} of odometry it would have to work with on hardware, from
the start, avoided deferring that discrepancy to the eventual
hardware migration.

Because Nav2, AMCL, and \texttt{bt\_navigator} all hard-code an
expectation that odometry arrives on \texttt{/odom} and driving
commands are consumed from \texttt{/cmd\_vel}, and the new controller
instead exposes these under its own namespaced topics
(\texttt{\textasciitilde/odometry}, \texttt{\textasciitilde/reference\_unstamped}),
two \texttt{topic\_tools relay} nodes were added to the launch files
to bridge the naming mismatch without modifying Nav2 itself.

\subsection{A Wheel-Axis Bug Exposed by the Migration}
\label{sec:wheel-axis-bug}

The migration exposed a latent URDF defect that had been silently
compensated for by the old manual node and had therefore never
produced a visible symptom before. The rear-left drive wheel joint
(\texttt{base\_wheelrl\_joint}) had its \texttt{<origin>} yawed by
\ang{180} for correct mesh mirroring, but retained the \emph{same}
\texttt{<axis>} direction as the corresponding right-side joint --
which, combined with the yaw flip, silently inverted that wheel's
effective rotation sense relative to a positive commanded velocity.
The old \texttt{cmd\_vel\_to\_drive.py} had manually compensated for
this by commanding \texttt{[-speed, +speed]} to the left/right wheel
velocities respectively; the new controller had no such
compensation, and computed wheel odometry that did not match the
vehicle's real motion direction. The fix -- correcting the
\texttt{<axis>} declaration on the affected joint directly, rather
than perpetuating a sign-flip workaround downstream in every
consumer of that joint -- illustrates a pattern worth stating
explicitly: a URDF geometry defect that is silently absorbed by one
consumer's ad hoc compensation will resurface, often with a confusing
symptom, the moment that consumer is replaced.

Listing~\ref{lst:wheel-axis} shows the two rear drive-wheel joints as
they stand in the corrected URDF: the left joint's mesh-mirroring yaw
(\texttt{rpy="0 0 3.14"}) is paired with a negated \texttt{<axis>} so
that a positive commanded wheel velocity turns both wheels in the
same physical sense, whereas the right joint needs no such
compensation since its origin carries no mirroring rotation.

\begin{lstlisting}[language=XML,caption={Rear-left and rear-right drive-wheel joints (\texttt{urdf/qcar\_model.xacro}) -- the axis on the mirrored left joint is negated to compensate for its \ang{180} origin yaw.},label={lst:wheel-axis}]
<joint name="base_wheelrl_joint" type="continuous">
    <parent link="base" />
    <child link="wheelrl" />
    <origin rpy="0 0 3.14" xyz="-0.12765 0.05610 0.03338" />
    <!-- axis negated to compensate for the 180deg yaw above (mesh mirror) -
         without this the left wheel spins opposite the right wheel for the
         same signed velocity command -->
    <axis xyz="0 -1 0" />
</joint>

<joint name="base_wheelrr_joint" type="continuous">
    <parent link="base" />
    <child link="wheelrr" />
    <origin rpy="0 0 0" xyz="-0.12765 -0.05610 0.03338" />
    <axis xyz="0 1 0" />
</joint>
\end{lstlisting}

\paragraph{Why the sign flip happens: axes live in the joint's own, rotated frame.}
The mechanism behind this bug is a direct consequence of one URDF
convention that is easy to overlook: a joint's \texttt{<axis>} vector
is expressed in the joint's \emph{own} local frame -- the frame
produced by applying that same joint's \texttt{<origin>} rotation to
its parent link's frame -- not directly in the parent link's (here,
\texttt{base}'s) frame. The rear-right joint's origin carries no
rotation (\texttt{rpy="0 0 0"}), so its joint frame is identical to
\texttt{base}'s frame, and \texttt{axis="0 1 0"} means exactly
$+Y$ in the vehicle's own reference frame. The rear-left joint's
origin, in contrast, is yawed \ang{180} about $Z$
(\texttt{rpy="0 0 3.14"}) so its mesh renders mirrored on the correct
side of the chassis. A \ang{180} yaw about $Z$ maps a frame's own $X$
and $Y$ axes to their negatives while leaving $Z$ unchanged
($X \to -X$, $Y \to -Y$), so that joint's \emph{local} $+Y$ direction
already points along \emph{base}'s $-Y$ direction before the
\texttt{<axis>} element is even considered. Declaring
\texttt{axis="0 1 0"} on the rear-left joint -- the value that looks,
by simple inspection, like ``the same axis as the right wheel'' --
therefore actually commands rotation about \texttt{base}'s $-Y$, the
opposite physical sense from the right wheel's \texttt{base}-frame
$+Y$, for the identical signed joint velocity. Negating it to
\texttt{axis="0 -1 0"} (Listing~\ref{lst:wheel-axis}) cancels the
already-applied yaw, restoring both wheels to rotating about the same
\texttt{base}-frame axis for the same commanded sign -- the fix is
not ``pick whichever sign makes it drive forward'' by trial and
error, but a direct, derivable consequence of the joint's own origin
rotation, which is why it generalises immediately to any other
mirrored-mesh joint pair in this or a similar URDF.

\subsection{A URDF Parser Defect in \texttt{gazebo\_ros2\_control}}
\label{sec:colon-bug}

While adding explanatory comments to the URDF for the wheel-axis fix
above, \texttt{controller\_manager} began failing to start
altogether -- with all downstream symptoms that implies
(\texttt{joint\_state\_broadcaster} and the new controller never
spawning, \texttt{/odom} never appearing, and no \texttt{map} TF
being published either, since nothing in the chain that would
produce it ever started). The failure was silent: no controller
spawner logged an error explaining \emph{why} it timed out waiting
for \texttt{controller\_manager}.

Isolating the cause required a delta-debugging-style bisection of the
URDF file against the actual parsing code path, since multi-line
comments, apostrophes, and file length were all plausible-looking red
herrings that turned out not to matter. The root cause was found in
how \texttt{gazebo\_ros2\_control} starts its internally-managed
\texttt{controller\_manager} node: the plugin passes the entire
processed URDF as a single \texttt{--param
robot\_description:=<xml>} command-line-style argument, which
\texttt{rcl}'s parameter-override grammar parses looking for a
\texttt{node\_name:param\_name} pattern. That grammar treats
\emph{any} colon immediately followed by a space --
\texttt{": "} -- anywhere in the argument string as a potential
separator, including inside an XML comment that a human reader would
never consider part of the "real" content. A direct reproduction
using \texttt{rclpy}'s own argument parser confirmed the mechanism
precisely: \texttt{"hello: world"} fails to parse while
\texttt{"hello:world"} (no space) succeeds, and \texttt{package://}
URIs or \texttt{foo:=bar} remappings are unaffected since neither
places a bare space directly after a colon.

This is not specific to this project's URDF; it is a property of any
\texttt{gazebo\_ros2\_control}-based setup, and is recorded here as a
general finding: \textbf{a URDF or xacro file must never contain the
literal sequence colon-followed-by-space, including inside comments},
or \texttt{controller\_manager} fails to start with no diagnostic
pointing at the actual cause.

\section{Porting the Navigation Stack to ROS~2 Humble / Gazebo Classic}
\label{sec:humble-conversion}

A sibling package in this project's workspace had been authored
against ROS~2 Jazzy and the newer Gazebo Sim (Harmonic) simulator --
the \texttt{gz sim} binary family, \texttt{ros\_gz\_sim} /
\texttt{ros\_gz\_bridge} integration, and \texttt{gz-sim-*-system}
world plugins -- none of which are available on this project's
ROS~2 Humble / Ubuntu 22.04 / Gazebo Classic~11 host. Rather than
install a second, newer simulator stack alongside the one already
validated and working for the rest of this thesis's packages
(Section~\ref{sec:first-migration}), the package was converted to
Gazebo Classic, matching every other package in this project.
Table~\ref{tab:gz-sim-conversion} summarises the conversion, which
generalises directly to any future package encountered that was
authored against Gazebo Sim / Jazzy-era tutorials.

\begin{table}[htbp]
  \centering
  \caption{Gazebo Sim (Jazzy) to Gazebo Classic (Humble) translation table.}
  \label{tab:gz-sim-conversion}
  \small
  \begin{tabular}{@{}p{4.2cm}p{4.2cm}p{5.5cm}@{}}
    \toprule
    Gazebo Sim (source) & Gazebo Classic (target) & Notes \\
    \midrule
    \texttt{ros\_gz\_sim} / \texttt{ros\_gz\_bridge} (package.xml) & \texttt{gazebo\_ros} / \texttt{gazebo\_plugins} & \\
    \texttt{gpu\_lidar} sensor + gz-sim system plugin & \texttt{type="ray"} + \texttt{\seqsplit{libgazebo\_ros\_ray\_sensor.so}} & \\
    gz-sim \texttt{imu} sensor & classic \texttt{imu} + \texttt{\seqsplit{libgazebo\_ros\_imu\_sensor.so}} & \\
    \texttt{\seqsplit{gz-sim-joint-state-publisher-system}} & removed entirely & Not needed; the drive plugin's own \texttt{publish\_wheel\_tf} broadcasts wheel/steering TF directly \\
    \texttt{\seqsplit{gz-sim-ackermann-steering-system}} & \texttt{\seqsplit{libgazebo\_ros\_ackermann\_drive.so}} & Classic plugin auto-computes wheelbase/track/radius from URDF geometry; the gz-sim system required stating them explicitly \\
    5$\times$ \texttt{<plugin filename="gz-sim-*-system">} (world SDF) & deleted entirely & Classic's \texttt{gzserver} has physics, scene broadcast, and sensors built into its core, not opt-in systems \\
    \texttt{gz sim} + \texttt{ros\_gz\_sim create} spawn + \texttt{ros\_gz\_bridge} node & \texttt{gazebo}/\texttt{gzserver} + \texttt{gazebo\_ros spawn\_entity.py} & No bridge node needed; classic plugins publish onto ROS~2 topics directly \\
    \texttt{GZ\_SIM\_RESOURCE\_PATH} env-var for mesh resolution & not needed & \texttt{gazebo\_ros} registers its own \texttt{package://} resolver \\
    \bottomrule
  \end{tabular}
\end{table}

One element of the original file required \emph{no} change: the
\texttt{<gazebo reference="link">\allowbreak<material>Gazebo/X</material>\allowbreak</gazebo>}
material-script blocks are classic-Gazebo-native syntax that happened
to already be present and compatible, a reminder that a systematic
conversion still benefits from checking each element rather than
assuming every gz-sim-authored construct needs translating.

\section{Second Migration: The Direct Drive Plugin}
\label{sec:second-migration}

The navigation stack's final drive architecture -- the one this
thesis's hardware bridge (Chapter~\ref{ch:hardware-bridge}) targets
-- drops \texttt{ros2\_control} entirely in favour of driving the
robot through \texttt{libgazebo\_ros\_ackermann\_drive.so}\allowbreak{} directly,
declared in a \texttt{<gazebo>} block in the URDF rather than through
a separate controller-manager configuration file, matching the
approach already adopted for the sibling package during the Humble
conversion above. This plugin subscribes to \texttt{/cmd\_vel} and
publishes \texttt{/odom} and the \texttt{odom} $\rightarrow$
\texttt{base} TF transform directly, again auto-computing wheelbase,
track width, and wheel radius from URDF joint and collision geometry
at spawn time -- removing an entire category of configuration
drift risk where a controller's stated geometry parameters could
silently fall out of sync with the URDF's actual geometry, which had
been possible (if not observed) under the first migration's
\texttt{ros2\_control} YAML-based configuration. This is the
architecture described in Section~\ref{sec:gazebo}'s account of the
drive plugin, and the one whose real-world calibration
Chapter~\ref{ch:calibration} presents.

\subsection{A Quiet Regression: Odometry Reverts to Ground Truth}
\label{sec:odom-regression}

This migration has one consequence worth stating explicitly, because
it is easy to overlook and was not caught at the time: it silently
reversed the odometry-source decision made in
Section~\ref{sec:first-migration}. Reading
\texttt{gazebo\_ros\_ackermann\_drive}'s own
\texttt{UpdateOdometryWorld()} implementation directly shows that its
\texttt{/odom} is computed from \texttt{model\_->WorldPose()} and
\texttt{model\_->WorldLinearVel()} -- the physics engine's own exact,
noise-free simulated pose and velocity for the whole vehicle body --
not an integration of individual wheel-joint velocities or steering
angles at all. In other words, the final simulation-side drive
architecture publishes \texttt{/odom} that is, mechanically, no
different from what \texttt{gazebo\_ros\_p3d} (Section~\ref{sec:first-migration})
would have provided: perfect, driftless ground truth, despite carrying
the same topic name and message type the first migration deliberately
chose \emph{instead of} ground truth.

This is not a bug in the sense the rest of this chapter uses the
word -- the plugin behaves exactly as its own source dictates, and
nothing in this project's testing depended on \texttt{/odom} actually
drifting in simulation. It is, however, a genuine and previously
unrecorded discrepancy between this thesis's stated intent (train
against the same \emph{kind} of odometry the physical robot will
have) and what the final simulation configuration actually exercises,
and it is worth stating plainly rather than leaving the earlier
claim uncorrected: \textbf{only the physical QCar's own odometry
(Chapter~\ref{ch:hardware-bridge}, Section~\ref{sec:onboard-bridge-odom})
is genuine encoder-based dead reckoning in this project; the
simulation's final \texttt{/odom} is ground truth, exactly like the
\texttt{gazebo\_ros\_p3d} source it was originally meant to replace}.
Framed in terms of Section~\ref{sec:sim-to-real-concept}'s sim-to-real
gap taxonomy, this means the localisation/odometry accuracy gap
between simulation and hardware in this project is not purely an
empirically discovered discrepancy -- it is, in part, structurally
guaranteed by this architecture, since the simulated AMCL/Nav2 stack
is given a perfect motion estimate that its real-hardware counterpart
can never have. Any future comparison of localisation accuracy
between simulation and hardware should account for this rather than
treating the two \texttt{/odom} streams as like for like.

\section{The ``Robot Not Moving Forward'' Bug Pattern}
\label{sec:not-moving-bug}

Both the sibling package's Humble conversion (Section~\ref{sec:humble-conversion})
and this project's own second migration (Section~\ref{sec:second-migration})
independently produced the same surface symptom immediately
afterward: the simulated vehicle did not drive forward under a
sustained \texttt{/cmd\_vel} command. In both cases, this single
symptom turned out to be caused by a stack of independent bugs, each
one masking the next until the previous was fixed -- a pattern
significant enough, and confirmed to recur closely enough between
the two packages, to warrant treating as a general defect signature
for this plugin/vehicle combination rather than two unrelated
incidents. Table~\ref{tab:not-moving-bugs} summarises all bugs found
across both occurrences.

\begin{table}[htbp]
  \centering
  \caption{The recurring ``not moving forward'' bug pattern, as found in both affected packages.}
  \label{tab:not-moving-bugs}
  \small
  \begin{tabular}{@{}p{0.6cm}p{5.3cm}p{7cm}@{}}
    \toprule
    \# & Bug & Symptom / mechanism \\
    \midrule
    1 & Wheel \texttt{<collision>} was the visual STL mesh, not a primitive shape & The plugin's radius auto-detection only supports cylinder/sphere collisions; a mesh collision silently returns radius 0, so \texttt{target\_linear / wheel\_radius} becomes infinite, poisoning the velocity PID regardless of gain. Fixed by keeping the STL for \texttt{<visual>} but using an explicit \texttt{<cylinder radius="0.036" length="0.0245"/>} for \texttt{<collision>}. \\
    2 & Zero Coulomb friction on the wheel links & Even with a correct wheel radius, the wheel spins at exactly the commanded angular velocity forever while the chassis never translates -- 100\% slip. Fixed with \texttt{<mu1>1.0</mu1><mu2>1.0</mu2>} on every wheel link. \\
    3 & Absurd ground-plane friction (\texttt{mu=100}, \texttt{mu2=50}) & Once wheel friction is fixed, this combination produces a real physics instability/NaN blow-up on contact. Fixed by lowering to a realistic \texttt{mu=1}, \texttt{mu2=1}. \\
    4 & Razor-thin wheel/ground clearance (only $\sim$0.38\,mm geometric overlap at the original 0.033\,m collision radius) & Enough overlap to rest on the ground without falling through, but not enough consistent contact-normal force to generate real traction even at correct friction -- chassis twist stayed near zero despite full wheel spin. Fixed by padding the collision cylinder radius to 0.036\,m (+3\,mm), which alone flipped the vehicle from ``spins in place'' to sustained forward motion. \\
    \bottomrule
  \end{tabular}
\end{table}

Listing~\ref{lst:wheel-collision-friction} shows the corrected form of
one rear wheel link and its friction declaration, combining the fixes
for bugs~1 and~2 from Table~\ref{tab:not-moving-bugs}: an explicit
cylinder \texttt{<collision>} geometry (the \texttt{<visual>} mesh,
omitted here, is unaffected and keeps using the STL), and a separate
\texttt{<gazebo reference="...">} block assigning Coulomb friction
coefficients to that same link -- URDF has no wheel-friction concept
of its own, so this Gazebo-specific extension is the only place that
value is set at all.

\begin{lstlisting}[language=XML,caption={Wheel collision cylinder and friction (\texttt{urdf/qcar\_model.xacro}) -- fixes for bugs 1 and 2 in Table~\ref{tab:not-moving-bugs}.},label={lst:wheel-collision-friction}]
<link name="wheelrl">
    <collision>
        <origin rpy="1.5708 0 0" xyz="0 0 0" />
        <geometry>
            <cylinder radius="0.036" length="0.0245" />
        </geometry>
    </collision>
    <!-- <visual> (STL mesh) and <inertial> omitted here, unaffected -->
</link>

<!-- wheel-ground friction - without this the wheels spin freely with 100%
     slip and the chassis never translates, regardless of PID gains -->
<gazebo reference="wheelrl">
    <mu1>1.0</mu1>
    <mu2>1.0</mu2>
</gazebo>
\end{lstlisting}

\paragraph{Mechanism: how bugs 1 and 2 poison the control loop and physics, in numbers.}
Section~\ref{sec:gazebo-pid-loop} and Section~\ref{sec:friction-theory}
give the general theory bugs~1 and~2 exercise; concretely, for this
vehicle, a commanded \SI{0.3}{\meter\per\second} with a mesh-shaped
(bug~1, radius $\to 0$) rear-wheel collision computes a target angular
velocity of $0.3 / 0 \to \infty$ rad/s in
Equation~\ref{eq:pid}'s error term -- not a large but finite number
that a differently-tuned PID could still track, but a value with no
finite representation at all, driving the linear-velocity PID's
output straight to its saturation limit on the very first physics
step regardless of gain. With the corrected \SI{0.036}{\meter} radius
(Listing~\ref{lst:wheel-collision-friction}) that same command instead
computes a finite target of $0.3 / 0.036 \approx \SI{8.33}{\radian\per\second}$,
a well-posed, trackable error signal. Bug~2 (zero friction) then
governs whether that correctly-spinning wheel produces any chassis
motion at all: with $\mu_1 = \mu_2 = 0$, Section~\ref{sec:friction-theory}'s
Coulomb limit permits zero tangential contact force regardless of how
much normal force the vehicle's weight presses the wheel into the
ground with, so the wheel can spin at exactly the commanded
\SI{8.33}{\radian\per\second} indefinitely while contributing zero
net force to the chassis -- the two bugs therefore fail at two
different, independently necessary conditions (a well-posed control
error, and a nonzero force limit to actually convert correct wheel
spin into chassis thrust), which is why fixing either one alone still
left the vehicle stationary and only fixing both together produced
motion.

The first occurrence of this pattern also included a preceding, more
basic defect specific to that package: its mesh URIs
(\texttt{package://.../models/qcar/QCarBody.stl}-style paths) pointed
at a directory layout copy-pasted from this project's own package
rather than the sibling's actual \texttt{meshes/} layout and file
names, so the robot had \emph{no} collision geometry at all and fell
through the simulated ground plane -- confirmed directly via
\texttt{/odom} tracking a $z$-position reaching approximately
$-38{,}000$\,m. This is worth recording as a separate, more
fundamental lesson: Gazebo Classic's URDF-to-SDF loader silently
rewrites \texttt{package://} mesh URIs to \texttt{model://}, which
only resolves via the \texttt{GAZEBO\_MODEL\_PATH} environment
variable -- and since neither package's \texttt{package.xml} exports
the \texttt{<gazebo\_ros gazebo\_model\_path=.../>} tag required for
that resolution, correcting the path string alone would not have
been sufficient. The fix instead used xacro's \texttt{\$(find)}
substitution to build \texttt{file://} URIs directly, which bypass
package-path resolution entirely.

The second occurrence additionally required a numerical-stability
fix not initially anticipated: wheel rotational inertia values that
were physically correct for a solid disk of the wheel's real mass and
radius, but small enough in absolute terms to interact badly with the
plugin's 100\,Hz PID control loop, producing an
\texttt{Ogre::AxisAlignedBox::setExtents} assertion crash a few
seconds into driving. Bumping wheel inertia roughly 20$\times$
(\texttt{ixx=izz=0.0012}, \texttt{iyy=0.0022}) together with lowering
the steering and linear-velocity PID gains resolved the crash; this
value combination, once found working in the first occurrence, was
reused directly in the second and confirmed stable there as well.

A further gain-tuning finding, specific to the second occurrence and
worth stating as a general caution for reusing PID values across
otherwise-identical vehicle configurations: a steering PID gain
verified stable and accurate on a \emph{frictionless} wheel
configuration produced a real turning radius roughly $6\times$ larger
than commanded once realistic wheel friction (bug~2 above) was added,
because the same low gain that was previously sufficient (with no
resistive torque to overcome) could no longer turn the front hubs
against a now-gripping, rolling tyre. The fix -- raising
\texttt{left/right\_steering\_pid\_gain} from \texttt{0.02} to
\texttt{3.0} -- was verified directly via TF lookups on the
steering-hub joints during a commanded 1.0\,m-radius turn, confirming
hub angles converging to the correct $\sim$\ang{15}--\ang{18}
Ackermann inner/outer split, and via a circle fit to the resulting
odometry trace, which recovered a turning radius of 0.87\,m against
the 1.0\,m commanded -- a reasonable margin given the real-vs-ideal
gap between the single-track model of
Section~\ref{sec:ackermann-kinematics} and the vehicle's true
four-wheel geometry.

\subsection{Diagnostic Method}
\label{sec:diagnostic-method}

A methodological point generalises from both occurrences of this bug
pattern: none of these four bugs are visible from \texttt{ros2 topic
list} or a static inspection of the URDF alone. A simulated robot
that appears correctly configured can simultaneously be in freefall,
frozen inside a crashed-but-still-running \texttt{gzserver} process,
or driving its wheels at 100\% slip with zero real chassis motion --
and none of these states are distinguishable without directly
comparing \texttt{/odom} position over several seconds of sustained
commanded motion, cross-checked against the Gazebo server log for
silent \texttt{[Err]}/assertion lines. This became this project's
standard verification method for any subsequent URDF or physics
change to the simulation stack.

\section{A Behaviour-Tree Fix for Ackermann Recovery}
\label{sec:bt-fix}

\paragraph{How a behaviour tree actually executes.}
A behaviour tree is evaluated by repeatedly ``ticking'' it from the
root: each tick propagates down the tree and every node it reaches
returns one of \texttt{SUCCESS}, \texttt{FAILURE}, or
\texttt{RUNNING} (still in progress, re-tick next cycle) back up to
its parent, and how a parent node interprets its children's results
depends entirely on that parent's type. A \texttt{RoundRobin} node --
the one this section's fix concerns -- ticks exactly one of its
children per activation (advancing to the next child, wrapping back
to the first, only once the currently-selected child itself returns
\texttt{SUCCESS} or \texttt{FAILURE} rather than \texttt{RUNNING}),
cycling through its full child list one at a time across repeated
recovery attempts rather than trying all of them on a single
activation. A \texttt{ReactiveFallback} (the parent of
\texttt{RecoveryActions} in Listing~\ref{lst:bt-recovery}) instead
re-ticks its children from the first one every single tick, moving on
to the next only if the current one fails -- the two control-flow
types are not interchangeable, and Nav2's stock recovery subtree pairs
them deliberately: the outer \texttt{ReactiveFallback} re-checks
whether a goal update has arrived (via the \texttt{GoalUpdated}
condition node) before committing to whichever recovery action the
inner \texttt{RoundRobin} has advanced to.

A separate, functional (rather than purely physics-level) issue
appeared once the vehicle was reliably driving: Nav2's default
recovery behaviour tree round-robins through
\texttt{[ClearCostmaps, Spin(1.57\,rad), Wait(5\,s), BackUp(0.3\,m)]}
whenever the active controller reports it cannot make further
progress. The \texttt{Spin} recovery commands a pure in-place
rotation -- physically impossible for an Ackermann-steered vehicle,
whose steering angle has no effect while linear velocity is zero, so
the vehicle simply cannot rotate without also translating. Roughly
one in three recovery attempts was therefore a guaranteed no-op that
burned its full timeout without attempting any corrective motion.

The fix was to remove \texttt{<Spin>} from the recovery action list
in local copies of Nav2's stock behaviour-tree XML files
(\texttt{navigate\_to\_pose\_w\_replanning\_and\_recovery.xml} and
its through-poses equivalent), wired into the launch configuration
via a \texttt{RewrittenYaml} substitution. This required one
non-obvious step: \texttt{RewrittenYaml} only overwrites
\emph{existing} leaf keys in a YAML file, it does not add new ones,
so the parameter keys pointing at these custom BT file paths
(\texttt{default\_nav\_to\_pose\_bt\_xml} /
\texttt{default\_nav\_through\_poses\_bt\_xml}) had to already exist
as placeholder entries in \texttt{nav2\_params.yaml} before the
launch-time substitution could take effect.

Listing~\ref{lst:bt-recovery} shows the resulting recovery fallback
from the local behaviour-tree copy: a \texttt{RoundRobin} of clearing
both costmaps, waiting, and backing up -- with no \texttt{<Spin>}
node present at all, compared to the stock tree's four-way rotation
through \texttt{[ClearCostmaps, Spin, Wait, BackUp]}.

\begin{lstlisting}[language=XML,caption={Recovery actions in the local behaviour-tree copy (\texttt{config/nav2/behavior\_trees/navigate\_to\_pose\_w\_replanning\_and\_recovery.xml}) -- no \texttt{<Spin>} node, unlike the Nav2 stock tree.},label={lst:bt-recovery}]
<ReactiveFallback name="RecoveryFallback">
  <GoalUpdated/>
  <RoundRobin name="RecoveryActions">
    <Sequence name="ClearingActions">
      <ClearEntireCostmap name="ClearLocalCostmap-Subtree"
        service_name="local_costmap/clear_entirely_local_costmap"/>
      <ClearEntireCostmap name="ClearGlobalCostmap-Subtree"
        service_name="global_costmap/clear_entirely_global_costmap"/>
    </Sequence>
    <Wait wait_duration="5"/>
    <BackUp backup_dist="0.30" backup_speed="0.05"/>
  </RoundRobin>
</ReactiveFallback>
\end{lstlisting}

Because \texttt{RoundRobin} advances exactly one position per
activation, dropping \texttt{Spin} from a four-child list to this
three-child list is not merely ``the useless action happens less
often'' -- it removes the guaranteed no-op activation entirely, so
every single recovery activation now attempts a physically achievable
action, and the cycle back to any one specific action (e.g.\ back to
\texttt{BackUp} after it was last tried) also completes one activation
sooner than it did with four children in the rotation.

\section{Summary of the Simulation-Stack Architecture}
\label{sec:sim-stack-summary}

By the end of this chapter's work, the simulation-side navigation
stack consists of: a URDF-described vehicle driven directly by
\texttt{libgazebo\_ros\_ackermann\_drive.so} (Section~\ref{sec:second-migration}),
publishing ground-truth-derived \texttt{/odom} (Section~\ref{sec:odom-regression}
-- a quiet regression from the genuine dead-reckoning odometry the
first migration, Section~\ref{sec:first-migration}, had deliberately
adopted), with corrected wheel collision geometry, friction, and inertia
(Section~\ref{sec:not-moving-bug}), tuned PID gains for both drive
and steering, and a custom Ackermann-aware recovery behaviour tree
(Section~\ref{sec:bt-fix}). This is the configuration
Chapter~\ref{ch:hardware-bridge} treats as the reference to be
matched -- as closely as the physical robot's own constraints allow
-- when bridging to hardware.
